\documentclass[a4paper,11pt]{article}
% ---------- Кириллица: работает и в pdfLaTeX, и в XeLaTeX/LuaLaTeX ----------
\usepackage{iftex}
\ifPDFTeX
  \usepackage[T2A]{fontenc}
  \usepackage[utf8]{inputenc}
  \usepackage[russian]{babel}
\else
  \usepackage{fontspec}
  \IfFontExistsTF{CMU Serif}{\setmainfont{CMU Serif}}{%
    \IfFontExistsTF{Times New Roman}{\setmainfont{Times New Roman}}{\setmainfont{FreeSerif}}}
  \usepackage{polyglossia}
  \setdefaultlanguage{russian}
\fi
\usepackage{amsmath,amssymb,mathtools}
\usepackage[a4paper,margin=2cm]{geometry}
\usepackage[dvipsnames]{xcolor}
\usepackage{tikz}
\usetikzlibrary{arrows.meta,calc,patterns,decorations.pathreplacing,positioning,babel}
\usepackage[most]{tcolorbox}
\usepackage{enumitem}
\usepackage{titlesec}
\usepackage{array}
\usepackage{tabularx}
\usepackage{needspace}
\usepackage{cancel}
\usepackage[unicode,hidelinks]{hyperref}

\emergencystretch=3em
\tolerance=1500
\setlength{\parindent}{0pt}
\setlength{\parskip}{4pt}
\setlist{nosep,leftmargin=*,topsep=2pt}
\setlength{\tabcolsep}{4pt}

% ---------- оформление (как в angem.tex) ----------
\definecolor{secred}{RGB}{190,40,45}
\definecolor{subteal}{RGB}{0,110,120}
\definecolor{leadcol}{RGB}{0,120,110}
\titleformat{\section}{\Large\bfseries\color{secred}}{\thesection.}{0.5em}{}
\titleformat{\subsection}{\large\bfseries\color{subteal}}{\thesubsection.}{0.5em}{}
\titleformat{\subsubsection}{\normalsize\bfseries\color{subteal}}{\thesubsubsection.}{0.5em}{}
\titlespacing*{\section}{0pt}{14pt}{6pt}
\titlespacing*{\subsection}{0pt}{10pt}{4pt}
\titlespacing*{\subsubsection}{0pt}{8pt}{3pt}

\newtcolorbox{opr}{breakable,enhanced,colback=blue!4,colframe=blue!55!black,
  boxrule=0pt,leftrule=2.2pt,arc=0pt,left=6pt,right=4pt,top=3pt,bottom=3pt,
  before skip=6pt,after skip=6pt}
\newtcolorbox{teor}[1]{breakable,enhanced,colback=orange!7,colframe=orange!70!black,
  boxrule=0pt,leftrule=2.2pt,arc=0pt,left=6pt,right=4pt,top=3pt,bottom=3pt,
  before skip=6pt,after skip=6pt,before upper={\textbf{Теорема #1.}\ }}
\newtcolorbox{hwbox}{breakable,enhanced,colback=gray!6,colframe=gray!60,
  boxrule=0.4pt,arc=1pt,left=6pt,right=6pt,top=3pt,bottom=3pt,fontupper=\small}

\newcommand{\lead}[1]{\textbf{\color{leadcol}#1}\ }
\newcommand{\q}{\textbf{\color{red!70!black}?}\ }
\newcommand{\imp}{\textbf{\color{red!70!black}!}\ }
\newcommand{\zam}{\textbf{Замечание.}\ }
\newcommand{\prim}{\textbf{Пример.}\ }
\newcommand{\dok}{\textbf{Доказательство.}\ }
\renewcommand{\thefootnote}{\textcolor{red!75!black}{\arabic{footnote}}}
\newcommand{\qedr}{\hfill$\blacksquare$}


\begin{document}
\begin{center}
{\LARGE\bfseries\color{secred}Теоретическая информатика}\\[3pt]
{\large Конспект лекций}
\end{center}
\vspace{-2mm}

\tableofcontents
\newpage


\begin{center}\textbf{Основы теории информации}\end{center}

% =====================================================================
\section{Информация и данные}
% =====================================================================

\begin{opr}
\textbf{Определение.} \emph{Информация}~--- это совокупность сведений, полученных путём наблюдения
и изучения, расширяющих представление об объектах, явлениях окружающей среды,
их свойствах, состоянии и взаимосвязях.
\end{opr}

\begin{opr}
\textbf{Определение.} \emph{Данные}~--- это информация в форме, пригодной для её хранения,
обработки, передачи.
\end{opr}

Основоположник теории информации~--- \textbf{Шеннон}.

\subsection{Система передачи информации}

Обмен информацией происходит между людьми, между людьми и окружающим миром.
Информацию передаёт какой-либо объект, система.

\begin{center}
\resizebox{\textwidth}{!}{%
\begin{tikzpicture}[
  node distance=2.2cm,
  box/.style={draw, minimum height=1.1cm, minimum width=2cm, align=center},
  lbl/.style={font=\footnotesize, align=center},
  >=stealth]
  \node[box] (src) {Источник\\информации};
  \node[box, right=of src] (tx) {Передатчик};
  \node[box, right=of tx] (ch) {Канал\\связи};
  \node[box, right=of ch] (rx) {Приёмник};
  \node[box, right=of rx] (dst) {Потребитель\\информации};
  \node[box, below=1.3cm of ch] (noise) {Источник\\помех};
  \draw[->] (src) -- node[lbl, above] {сообщение} (tx);
  \draw[->] (tx) -- node[lbl, above] {передаваемый\\сигнал} (ch);
  \draw[->] (ch) -- node[lbl, above] {принятый\\сигнал} (rx);
  \draw[->] (rx) -- node[lbl, above] {сообщение} (dst);
  \draw[->] (noise) -- node[lbl, right] {помехи} (ch);
\end{tikzpicture}}
\end{center}

\textbf{Источник} отправляет информацию в виде сообщения.

\begin{opr}
\textbf{Определение.} \emph{Сообщение}~--- это форма представления информации для её
непосредственной передачи в виде данных.
\end{opr}
Формы сообщений:
% ниже список в оригинале читается лишь частично; восстановлен по смыслу
\begin{itemize}
  \item текстовая~--- текст на естественном языке;
  \item кодовая~--- закодированные данные;
  \item графическая~--- изображения объектов;
  \item звуковая;
  \item видео;
  \item сигналы.
\end{itemize}

\textbf{Передатчик} преобразует сообщение в сигналы, пригодные для передачи
по каналу связи. Например, в мобильной телефонии передатчик преобразует
звуковое сообщение в радиосигналы.

\begin{opr}
\textbf{Определение.} \emph{Сигнал} (в общем случае)~--- это физический динамический процесс,
при котором его параметры изменяются во времени.
\end{opr}

\begin{opr}
\textbf{Определение.} \emph{Канал связи}~--- это физическая среда, используемая для передачи
сигнала от передатчика к приёмнику. Каналом может быть коаксиальный кабель,
полоса радиочастот, луч света и т.\,д.
\end{opr}

\textbf{Приёмник} обычно выполняет преобразование, обратное преобразованию,
проводимому передатчиком. Восстанавливает сообщение из сигнала.

\textbf{Потребитель информации}~--- это человек или устройство, для которого
предназначено сообщение.

Чаще всего при передаче сигнала эти сигналы могут быть искажены помехами.
\textbf{Источник помех} действует на передаваемые сигналы, изменяя
принимаемое сообщение.

Система передачи информации проявляется в различных сферах нашей жизни:

\begin{center}
\small
\begin{tabularx}{\textwidth}{|X|X|X|X|}
\hline
\textbf{Части} & \textbf{Лекция} & \textbf{Запрос веб-сервера (ВС)} & \textbf{Телевидение} \\
\hline
Источник информации & Лектор & Клиентский компьютер (КК) & Источник радиосигнала \\
\hline
Передатчик & Голосовые связки лектора & Сетевая плата КК & Радиоканал \\
\hline
Канал связи & Воздушная среда & Провайдер, канал связи (ПКС) & Воздух, радиоканал \\
\hline
Приёмник & Уши & Сетевая плата & Антенна радиоприёмника \\
\hline
Потребитель информации & Студент & ВС & Слушатель \\
\hline
Источник помех & Шум в аудитории, уличный шум & Воздействие на ПКС & Атмосферные шумы \\
\hline
Сообщение & Текст & Текстовое или числовое & Видеоизображение \\
\hline
\end{tabularx}
\end{center}

% =====================================================================
\section{Дискретные и непрерывные сигнальные каналы}
% =====================================================================

\begin{opr}
\textbf{Определение.} Если параметр сигнала принимает конечное число значений и при этом они все
могут быть пронумерованы, то этот сигнал называется \emph{дискретным}.
Канал для передачи таких сигналов также называется дискретным.
\end{opr}

Примером дискретных сигналов являются буквы. Количество букв конечно, и их
можно рассматривать как уровни сигнала передачи сообщения.

\begin{opr}
\textbf{Определение.} Если параметр сигнала является непрерывной во времени функцией, то сигнал
называется \emph{непрерывным}. Канал для передачи непрерывных сигналов называется непрерывным.
\end{opr}

\prim Человеческая речь~--- передаваемая звуковой волной,
характеризуемой частотой, фазой и амплитудой. Параметром сигнала в этом
случае является давление, создаваемое этой волной в точке нахождения
приёмника~--- человеческого уха.

Для передачи информации используются непрерывные и дискретные каналы связи.
Сигналы переводят из непрерывных в дискретные и обратно.

\subsection{Дискретизация}

\begin{opr}
\textbf{Определение.} \emph{Дискретизация}~--- это процесс преобразования непрерывного сигнала
в дискретный с конечным числом значений на определённом отрезке.
\end{opr}

Для этого диапазон значений функции (ось ординат) разбивается на конечное
число отрезков равной ширины, называемой \emph{шагом дискретизации}.
Дискретное значение определяется отрезком, в который попало непрерывное значение функции.

\imp Чем меньше шаг дискретизации, тем ближе полученный дискретный сигнал к исходному
непрерывному сигналу, следовательно, больше степень дискретизации.

\prim На метеостанции каждые полчаса производят замер температуры.

\begin{center}
\begin{tikzpicture}[font=\small, >=stealth, x=1.1cm, y=0.55cm]
  \foreach \y in {1,...,5} \draw[gray!50] (0,\y) -- (7.6,\y);
  \draw[->] (-0.2,0) -- (7.8,0) node[below] {$t$, ч};
  \draw[->] (0,-0.2) -- (0,6.3) node[left] {$T$};
  \foreach \y in {1,...,5} \node[left] at (0,\y) {\footnotesize\y};
  \draw[thick] plot[smooth] coordinates {(0,3.0) (0.8,3.7) (1.6,4.3) (2.4,4.7) (3.1,4.75) (3.9,4.1) (4.8,3.0) (5.8,2.8) (7.4,3.8)};
  \foreach \x/\v/\l in {0.5/3/x_1, 1.0/4/x_2, 1.5/4/x_3, 2.0/5/x_4, 2.5/5/x_5, 3.0/5/x_6} {
    \draw[dashed, gray] (\x,0) -- (\x,\v);
    \fill[red!75!black] (\x,\v) circle (1.6pt);
    \node[below] at (\x,0) {\footnotesize$\l$};
  }
\end{tikzpicture}
\end{center}

Непрерывно меняющаяся температура замеряется в моменты $x_1, x_2, \dots$; в журнал записываются
значения, являющиеся дискретными значениями. В этом примере получасовой промежуток является шагом
дискретизации по времени, шаг дискретизации по значению равен 1, т.\,к. производится округление
до целого числа, а получасовые замеры температуры являются процессом дискретизации.

% =====================================================================
\section{Мера информации}
% =====================================================================

\subsection{Вероятность событий}

\[
  P = \frac{m}{N}, \qquad 0 \le P \le 1; \qquad\qquad
  P_i = \frac{m_i}{N}, \qquad \sum_{i=1}^{n} P_i = 1.
\]

\subsection{Мера информации Хартли}

Пусть система находится в одном из $n$ состояний в каждый момент времени, причём каждое из
состояний равновероятно. Например, в качестве системы могут выступать опыты с подбрасыванием
монеты ($n = 2$) или бросанием игральной кости ($n = 6$). Система может хранить количество единиц
информации, исчисляемое количеством состояний, по формуле Хартли:
\[
  \boxed{H = \log_b n = \frac{\ln n}{\ln b}}
\]

Основание $b$ определяет единицу измерения информации. При $b = 2$, $n = 2$ имеем $H = 1$~---
количество информации минимально, поэтому в качестве единицы информации принимают два
равновероятных состояния системы, например <<ложь>>~--- <<истина>>.

\begin{opr}
\textbf{Определение.} Такая единица количества информации называется \emph{бит}
(\emph{bit}~--- \emph{binary digit}).
\end{opr}

\subsection{Энтропия}

\begin{opr}
\textbf{Определение.} \emph{Энтропия}~--- мера информации Шеннона (рассматривается среда передачи
информации с дискретным каналом без помех). Пусть в сообщениях используются символы алфавита $A$
с вероятностями $(p_1, p_2, \dots, p_n)$. Мера информации, которую в среднем несёт символ алфавита,
называется \emph{энтропией} $H$ этого алфавита. Энтропия определяется формулой
\[
  H = \sum_{i=1}^{n} p_i \log_2 \frac{1}{p_i} = -\sum_{i=1}^{n} p_i \log_2 p_i.
\]
\end{opr}

Смысл энтропии может быть рассмотрен в различных аспектах:
\begin{enumerate}[label=\arabic*)]
  \item неопределённость для получателя информации, какой символ будет передан, о чём будет
        сообщение;
  \item объём информации, создаваемый источником, или среднее количество бит, требуемое для
        кодирования алфавита $A$.
\end{enumerate}

\lead{Свойства энтропии:}
\begin{enumerate}[label=\arabic*)]
  \item энтропия измеряется в битах на символ;
  \item энтропия максимальна, если состояния системы равновероятны;
  \item если вероятность одного символа равна единице, а остальных~--- 0, то энтропия равна 0:
        получатель заранее знает получаемые символы, и неопределённость равна 0.
\end{enumerate}

\lead{Применения энтропии:}
\begin{enumerate}[label=\arabic*)]
\item \textbf{Основная теорема Шеннона для канала без шума.}
\begin{teor}{(Шеннона)}
  Пусть источник имеет энтропию $H$ (бит на символ), а канал имеет пропускную способность $C$
  (бит в секунду). Тогда можно закодировать сообщение таким образом, чтобы передавать символы
  по каналу со средней скоростью
  \[
    \frac{C}{H} - \varepsilon \quad \text{(символов в секунду)},
  \]
  где $\varepsilon$ сколь угодно мало. Со средней скоростью больше, чем $\dfrac{C}{H}$, невозможно.
\end{teor}
Для приближения к этому пределу необходимо более частые символы представлять более короткими кодами.

\item \textbf{Относительная энтропия.}
\begin{opr}
\textbf{Определение.} Отношение энтропии к максимальному значению, которое могла бы достичь
энтропия при тех же символах, называется \emph{относительной энтропией}.
\end{opr}
Относительная энтропия является мерой максимального сжатия информации, которое достигается
кодированием при помощи того же самого алфавита. Единица минус относительная энтропия есть
\emph{избыточность}. Избыточность обычного английского текста составляет примерно 50\%.

\item \textbf{Энтропия определяет степень сжатия символьными кодами.} Символьные коды представляют
исходный алфавит так, чтобы обеспечить максимальное сжатие. Пусть элементы сообщения кодируются
кодами длины $(l_1, l_2, \dots, l_n)$. \emph{Ожидаемая длина кода} определяется по формуле
\[
  L = \sum_{i=1}^{n} p_i l_i.
\]
Ожидаемая длина кода ограничена энтропией: $\boxed{H \le L < H + 1}$.
\end{enumerate}

\prim В сообщении используются независимые буквы $A, B, C, D$ с вероятностями
$\frac12, \frac14, \frac18, \frac18$.
\[
  H = -\left(\tfrac12 \log_2 \tfrac12 + \tfrac14 \log_2 \tfrac14 + \tfrac18 \log_2 \tfrac18 + \tfrac18 \log_2 \tfrac18\right)
    = \frac74 = 1{,}75 \ \text{бит на символ}.
\]
Максимальная энтропия для букв $A, B, C, D$: $H_{\max} = \log_2 4 = 2$ при их вероятностях
$\frac14, \frac14, \frac14, \frac14$.

Относительная энтропия: $\dfrac{H}{H_{\max}} = \dfrac{1{,}75}{2} = 0{,}875$, \ а избыточность:
$1 - \dfrac{H}{H_{\max}} = 0{,}125$.

Зададим коды букв двумя способами:
\begin{enumerate}[label=\arabic*)]
  \item с учётом их вероятности (более частые буквы кодируются более короткими кодами);
  \item без учёта.
\end{enumerate}

\begin{center}
\renewcommand{\arraystretch}{1.2}
\begin{tabular}{c|c|c}
  Буква & Способ 1 (с учётом) & Способ 2 (без учёта) \\ \hline
  $A$ & 0   & 00 \\
  $B$ & 10  & 01 \\
  $C$ & 111 & 10 \\
  $D$ & 110 & 11 \\
\end{tabular}
\end{center}
Ожидаемая длина кода $L = \sum p_i l_i$:
\[
  L_1 = \tfrac12 \cdot 1 + \tfrac14 \cdot 2 + \tfrac18 \cdot 3 + \tfrac18 \cdot 3 = \tfrac74, \qquad
  L_2 = \tfrac12 \cdot 2 + \tfrac14 \cdot 2 + \tfrac18 \cdot 2 + \tfrac18 \cdot 2 = 2.
\]
Первый способ достигает нижней границы: $L_1 = H$.

\prim Буквы $A, B, C, D, E$ с вероятностями $0{,}3$; $0{,}35$; $0{,}1$; $0{,}1$; $0{,}15$.
Коды строятся по дереву: на каждом ветвлении левой ветви приписывается 0, правой~--- 1.

\begin{center}
\begin{tikzpicture}[font=\small, level distance=0.9cm,
  level 1/.style={sibling distance=3.2cm}, level 2/.style={sibling distance=1.6cm},
  level 3/.style={sibling distance=1.1cm},
  every node/.style={inner sep=1.5pt}]
  \node[circle,fill,inner sep=1.3pt] {}
    child { node[circle,fill,inner sep=1.3pt] {}
      child { node {$A$} edge from parent node[left] {0} }
      child { node[circle,fill,inner sep=1.3pt] {}
        child { node {$C$} edge from parent node[left] {0} }
        child { node {$D$} edge from parent node[right] {1} }
        edge from parent node[right] {1} }
      edge from parent node[left] {0} }
    child { node[circle,fill,inner sep=1.3pt] {}
      child { node {$B$} edge from parent node[left] {0} }
      child { node {$E$} edge from parent node[right] {1} }
      edge from parent node[right] {1} };
\end{tikzpicture}
\qquad
\renewcommand{\arraystretch}{1.15}
\begin{tabular}[b]{c|c|c}
  Буква & $p_i$ & Код \\ \hline
  $A$ & 0,3  & 00 \\
  $B$ & 0,35 & 10 \\
  $C$ & 0,1  & 010 \\
  $D$ & 0,1  & 011 \\
  $E$ & 0,15 & 11 \\
\end{tabular}
\end{center}

\begin{hwbox}
\textbf{Самостоятельно:} объёмный подход (раздел 2.4.2), разделы 2.5 и 2.3.
\end{hwbox}

% =====================================================================
\section{Системы счисления}
% =====================================================================

\subsection{Число и его модель}

Числа занимают важную часть современного мира:
\begin{enumerate}[label=\arabic*)]
  \item результаты измерений обычно записываются числами;
  \item в ЭВМ (электронной вычислительной машине) обработка данных сведена к вычислениям.
        Исходными и конечными данными вычислений являются числа.
\end{enumerate}

В природе количество чисел не ограничено, однако практические модели содержат конечное
количество чисел. Например, в ЭВМ могут храниться очень большие числа, но их диапазон ограничен.
Наиболее точной моделью числа является его запись. Запись числа необходима для его запоминания.

\subsection{Система счисления. Определение и характеристики}

Числа записываются в системе счисления.

\begin{opr}
\textbf{Определение.} \emph{Система счисления}~--- это соглашение о представлении чисел посредством
конечной совокупности символов (цифр)
\[
  A = \{a_1, a_2, \dots, a_n\}.
\]
$A$ называется \emph{алфавитом}; каждой цифре ставится соответствие, определяющее её
количественный эквивалент.
\end{opr}

Люди используют десятичную систему счисления, а ЭВМ~--- двоичную (устройства с двумя
состояниями проще).

\lead{Характеристиками системы счисления являются:}
\begin{enumerate}[label=\arabic*)]
  \item отличимость символов;
  \item вычислимость чисел;
  \item выявляемость ошибок в записи.
\end{enumerate}

Одна и та же система счисления имеет разные показатели характеристик для человека и ЭВМ.
Например, вычисления в двоичной системе счисления для ЭВМ проще, чем для человека, а в десятичной~---
наоборот.

\subsection{Позиционные и непозиционные системы счисления}

Системы счисления разделяют на позиционные и непозиционные.

\begin{opr}
\textbf{Определение.} \emph{Непозиционные системы счисления}~--- это системы, в которых цифры
не меняют своего количественного эквивалента в зависимости от местоположения (позиции) в записи
числа.
\end{opr}
К непозиционным системам счисления относится система римских цифр.

\begin{opr}
\textbf{Определение.} \emph{Позиционные системы счисления}~--- это системы, в которых
количественный эквивалент цифры зависит от её положения в числе.
\end{opr}
Пример~--- десятичная система счисления.

\begin{opr}
\textbf{Определение.} \emph{Основание} позиционной системы счисления~--- это количество символов в её
алфавите. Позиционные системы счисления с основанием $n$ называют \emph{$n$-ичными}.
\end{opr}
В десятичной СС 10 цифр, поэтому она имеет основание $n = 10$.

\imp Преимущество позиционной СС: для записи числа $K$ необходимо не более $\bigl[\log_n K\bigr] + 1$
цифр. Операции над числами сводятся к действиям над его цифрами.

Далее рассматриваются только позиционные системы счисления, поэтому слово <<позиционная>> опускается.

\subsection{Позиционное и полиномиальное представление чисел}

\prim Полином: \ $a_0 x^0 + a_1 x^1 + a_2 x^2 + \dots + a_n x^n$, \ например $2x^3 + 4x + 5$.

\begin{opr}
\textbf{Определение.} \emph{Позиционное представление} числа $B_n$ с целой и дробной частью
в $n$-ичной СС имеет форму
\[
  B_n = a_m a_{m-1} \dots a_1 a_0,\, a_{-1} a_{-2} \dots a_{-k},
\]
запятая разделяет целую (слева) и дробную (справа) части; $a_i$, $i = -k, \dots, m$~--- цифра
в позиции $i$; $n$~--- явное указание основания системы счисления.
\end{opr}

\begin{opr}
\textbf{Определение.} \emph{Полиномиальное представление} $B_n$ имеет вид суммы произведений цифр
на вес их позиции в числе:
\[
  B_n = a_m n^m + a_{m-1} n^{m-1} + \dots + a_1 n^1 + a_0 n^0 + a_{-1} n^{-1} + a_{-2} n^{-2}
        + \dots + a_{-k} n^{-k} = \sum_{i=-k}^{m} a_i n^i. \tag{4.1}
\]
\end{opr}

\imp Позиционное представление определяет запись (вид числа), а полиномиальное~--- значение числа.

\subsection{Двоичная, десятичная, шестнадцатеричная системы счисления}

В информатике используются следующие системы счисления:
\begin{enumerate}[label=\arabic*)]
  \item двоичная: $n = 2$, алфавит $A = \{0, 1\}$;
  \item десятичная: $n = 10$, алфавит $A = \{0, 1, 2, 3, 4, 5, 6, 7, 8, 9\}$;
  \item шестнадцатеричная: $n = 16$, алфавит $A = \{0, 1, \dots, 9, A, B, C, D, E, F\}$.
\end{enumerate}

\begin{hwbox}
\textbf{Д/з:} DRAM-память (посмотреть в интернете).
\end{hwbox}

В вычислительной технике используется двоичная система счисления. Все числа и данные
представляются в виде последовательности $\{0, 1\}$ (бит). Двоичная СС обладает следующими
преимуществами перед СС с другими основаниями:
\begin{enumerate}[label=\arabic*)]
\item Для хранения двоичных цифр необходимы технические устройства с двумя устойчивыми
      состояниями, которые легко различимы. Например: ток есть~--- ток отсутствует,
      намагничено~--- не намагничено и т.\,п.
\item Представление данных посредством только двух состояний надёжно и помехоустойчиво.

\begin{center}
\begin{tikzpicture}[font=\small, >=stealth]
  \draw[thick] (0,0) rectangle (6,0.6);
  \draw[thick] (4.4,0) -- (4.4,0.6);
  \draw[thick] (4.9,0) -- (4.9,0.6);
  \node at (2.2,0.3) {данные};
  \node at (4.65,0.3) {0};
  \node at (5.45,0.3) {\dots};
  \node[above] at (3,0.6) {CRC};
  \draw[->] (4.65,-0.9) node[below] {бит чётности} -- (4.65,-0.05);
  \draw[->] (1.3,-0.9) node[below, align=center] {исходные\\передаваемые данные} -- (1.3,-0.05);
\end{tikzpicture}
\end{center}
Помехоустойчивость~--- к данным добавляются контрольные разряды (контрольная сумма CRC, бит
чётности). Если данные чётные и бит чётности чётный, то всё OK.

\item Для выполнения арифметических операций используется простой аппарат логики высказываний
      (булевой алгебры).
\end{enumerate}

Шестнадцатеричная система счисления используется для сокращения записи чисел в двоичной СС.

\subsection{Преобразование чисел из одной системы счисления в другую}

В вычислительной технике ввод, вывод и обработка числовых данных связаны с преобразованием чисел
из одной СС в другую.

\begin{opr}
\textbf{Правило 4.1} (преобразование целого или дробного числа из $n$-ичной СС в десятичную).
Число из $n$-ичной СС преобразуется в десятичную СС по формуле (4.1).
\end{opr}

\prim $11001_2 = 1 \cdot 2^4 + 1 \cdot 2^3 + 0 \cdot 2^2 + 0 \cdot 2^1 + 1 \cdot 2^0 = 16 + 8 + 1 = 25_{10}$.

При преобразовании из одной СС в другую (недесятичную) правила преобразования различаются для
целых и дробных чисел.

\begin{opr}
\textbf{Правило 4.2} (преобразование целых чисел из 10-ичной СС в $n$-ичную).
Целое число $X$ преобразуется в $n$-ичную СС по следующему алгоритму:
\begin{enumerate}[label=\arabic*)]
  \item получить цифру числа $n$-ичной СС как остаток деления числа $X$ на основание новой СС $n$;
        полученную цифру записать слева от предыдущих цифр;
  \item принять за $X$ частное от деления числа $X$ на основание СС $n$;
  \item выполнять пп.~1), 2), пока $X \ne 0$.
\end{enumerate}
\end{opr}

\textbf{Пример 4.3.} Перевести $25$ в СС с $n = 2$.
\[
  \begin{array}{rcll}
    25 : 2 &=& 12, & \text{остаток } 1 \\
    12 : 2 &=& 6,  & \text{остаток } 0 \\
     6 : 2 &=& 3,  & \text{остаток } 0 \\
     3 : 2 &=& 1,  & \text{остаток } 1 \\
     1 : 2 &=& 0,  & \text{остаток } 1
  \end{array}
  \qquad\Longrightarrow\qquad 25_{10} = 11001_2
\]
(остатки записываются справа налево: последний остаток~--- старшая цифра).

\begin{hwbox}
\textbf{Самостоятельно:} преобразование дробных чисел из одной СС в другую.
\end{hwbox}

% =====================================================================
\section{Представление чисел, символов, графических и звуковых данных в ЭВМ}
% =====================================================================

\subsection{Представление целых чисел в ЭВМ}

Целые числа~--- универсальный способ представления вещественных чисел, точек изображения и звука:
форматы представления вещественных чисел, точек изображения, звука используют целые числа как
базовый элемент.

\lead{Форматы представления целых чисел.} Числа в ЭВМ хранятся в соответствии с форматами.

\begin{opr}
\textbf{Определение.} \emph{Формат}~--- это соглашение или правило представления числа как
последовательности разрядов.
\end{opr}

Минимальная единица хранения данных в ЭВМ~--- 1 байт.

В настоящее время в 32-разрядных ЭВМ существуют следующие форматы представления целых чисел:
\emph{байт}, \emph{полуслово} (2 байта), \emph{слово} (4 байта), \emph{двойное слово} (8 байт).
Биты, составляющие эти форматы, называются \emph{разрядами}. Слева~--- старшие разряды,
справа~--- младшие.

Каждый из этих форматов может быть \emph{знаковым} для представления положительных или
отрицательных чисел, или \emph{беззнаковым} для представления положительных чисел.
Знаковым является самый старший разряд: если разряд $= 0$, то число положительное, если
разряд $= 1$, то число отрицательное.

В формате целых чисел размером больше 1 байта числа хранятся в памяти ЭВМ в обратном порядке:
вначале младший байт, потом старший.

\prim Число $\mathrm{B5DE}_{16}$ хранится как \ \fbox{DE}\fbox{B5}\,.

Такое расположение байт ускоряет операции с числами, т.\,к. вычисления по алгоритмам начинаются
с младших разрядов.

\subsection{Прямой и дополнительный коды представления двоичных чисел}

\begin{enumerate}[label=\arabic*)]
\item В \emph{прямом коде} старший бит кодирует знак числа (0 для <<$+$>>, 1 для <<$-$>>), а остальные
биты~--- модуль числа.

\prim (знаковый бит отделён чертой) \quad $11_{10} = 0|1011_{\text{п}}$, \qquad $-11_{10} = 1|1011_{\text{п}}$.

\item В \emph{дополнительном коде} положительные числа кодируются так же, как и в прямом. Для
преобразования отрицательного числа в дополнительный код есть 2 способа.
\end{enumerate}

\begin{opr}
\textbf{Определение.} Операция \emph{инвертирования}~--- замена значения разряда на противоположное.
\end{opr}

\begin{opr}
\textbf{Правило 5.1} (поразрядное представление числа в дополнительном коде).
Представить модуль отрицательного числа в прямом коде и проинвертировать все разряды левее самой
младшей (правой) единицы.
\end{opr}

\prim Представить число $-11_{10}$ в дополнительном коде:
\[
  |{-11}| = 0|1011 \quad\xrightarrow{\text{инвертируем всё левее младшей единицы}}\quad 10101 = -11_{10}.
\]

\begin{opr}
\textbf{Правило 5.2} (арифметическое представление отрицательного числа в дополнительном коде).
Прибавить к отрицательному числу $2^m$, где $m$~--- количество разрядов в двоичном представлении
или данном формате, и полученное число перевести в двоичную СС. Для байта $2^8 = 256$.
\end{opr}

\imp Из этих правил следует, что положительные числа в случае увеличения числа разрядов дополняются
слева нулями, а отрицательные~--- единицами.

\prim Представить $-11_{10}$ в дополнительном коде ($m = 5$):
\[
  2^5 = 32 \ \Rightarrow\ -11 + 32 = 21, \qquad 21_{10} = 10101_2.
\]

Взаимно обратное преобразование отрицательных чисел и чисел, записанных в дополнительном коде,
выполняется по тем же правилам.

\subsection{Операции с целыми числами}

Центральный процессор ЭВМ только складывает и выполняет операции с разрядами (логические~---
конъюнкция, дизъюнкция, инвертирование, сдвиг). Поэтому все арифметические операции сводятся
к сложению.

Запись в дополнительном коде используется для сложения и вычитания положительных и отрицательных
чисел без преобразования.

\prim Сложить $21_{10}$ и $-11_{10}$ (6 разрядов):
\[
  \begin{array}{r@{\;}l}
     & 010101 \quad (21_{10}) \\
    + & 110101 \quad (-11_{10}) \\ \hline
     & \cancel{1}001010 = 10_{10}
  \end{array}
\]
(перенос из старшего разряда отбрасывается).

При умножении чисел используется \emph{модифицированный алгоритм Бута}. Множители
представлены в дополнительном коде.

\begin{opr}
\textbf{Правило 5.6} (модифицированный алгоритм Бута). Числа умножаются следующим образом:
\begin{enumerate}[label=\arabic*)]
\item Разделить второй множитель справа налево на перекрывающиеся триады. Триады должны
      перекрываться в одном разряде. Если правая триада не полная, то дополнить её нулями.
\item Заменить триады на коды операций:
\begin{center}
\begin{tabular}{c|c}
  Триада & Код \\ \hline
  000 & $0$ \\
  001 & $+1$ \\
  010 & $+1$ \\
  011 & $+2$ \\
  100 & $-2$ \\
  101 & $-1$ \\
  110 & $-1$ \\
  111 & $0$
\end{tabular}
\end{center}
\item Для каждого полученного кода определить слагаемое как результат операции с первым множителем:
\begin{center}
\begin{tabular}{c|p{10cm}}
  Код & Операция с множителем \\ \hline
  $-2$ & получить дополнительный код и умножить на 2 (сдвиг на 1 бит влево) \\
  $-1$ & получить дополнительный код \\
  $0$  & сформировать строку из нулей \\
  $+1$ & не изменять \\
  $+2$ & умножить на 2 (сдвиг на 1 бит влево)
\end{tabular}
\end{center}
\item Сдвинуть каждое слагаемое влево на 2 разряда относительно предыдущего и сложить слагаемые.
      Для множителя из $k$ разрядов необходимо получить $\frac{k}{2}$ слагаемых.
\end{enumerate}
\end{opr}

\prim Умножить $-11_{10}$ на $29_{10}$.
\[
  -11_{10} = 10101_2 \ \text{(дополнительный код)}, \qquad
  29_{10} = 011101\underline{0}_2 \ \text{(дописан ноль справа)}.
\]
Триады справа налево: $010 \to +1$, \ $110 \to -1$, \ $011 \to +2$.

Слагаемые (10 разрядов, каждое следующее сдвинуто на 2 разряда влево):
\[
  \begin{array}{r@{\;}ll}
      & 1111110101   & (+1:\ -11) \\
    + & 0000101100   & (-1:\ +11,\ \text{сдвиг на 2}) \\
    + & 1010100000   & (+2:\ -22,\ \text{сдвиг на 4}) \\ \hline
      & 1011000001   & = -319_{10}
  \end{array}
\]
Проверка: модуль результата $0100111111_2 = 319_{10}$, и $-11 \cdot 29 = -319$.

При операциях с целыми числами возможно \emph{переполнение} или \emph{переполнение снизу}:
\begin{itemize}
  \item \emph{переполнение} (overflow) возникает при больших значениях представленного числа;
  \item \emph{переполнение снизу} (underflow) возникает при результате операции меньшем, чем
        представимое число.
\end{itemize}

\end{document}
